\section{Introduction}\label{Introduction}
Designing molecules requires balancing two distinct modalities: the molecular \emph{structure} of atoms and bonds, and the linguistic descriptions of molecular properties (e.g., ``a kinase inhibitor with good solubility''). Bridging this gap through cross-modal translation enables chemically grounded workflows where designs can be iteratively refined via natural language~\cite{zheng2023structureinformedlanguagemodelsprotein,fatemi2024talk}. This motivates two complementary directions: (i) \textbf{molecule generation} (\emph{text}$\rightarrow$\emph{molecule}) to materialize descriptions into structures, and (ii) \textbf{molecule captioning} (\emph{molecule}$\rightarrow$\emph{text}) to explain molecular structure and properties via natural-language descriptions~\cite{edwards-etal-2022-translation,zeng2023interactivemoleculardiscoverynatural}. However, despite their symmetry, these tasks are typically developed in isolation with separate models and objectives, limiting consistency and reuse in iterative design loops.

\noindent A common bridge between molecular structure and sequence models is the Simplified Molecular-Input Line-Entry System (SMILES)~\cite{doi:10.1021/ci00057a005,doi:10.1021/ci00062a008}, which linearizes a molecular graph into a string. This linearization has enabled a large family of sequence-based molecular methods, including (1) conditional SMILES generation from text prompts~\cite{doi:10.1021/acs.jcim.1c00600,D1RA03086H,irwin2022chemformer}, and (2) SMILES or structure-conditioned captioning that maps molecules to natural-language descriptions~\cite{edwards-etal-2022-translation,zeng2023interactivemoleculardiscoverynatural}. Many of these systems rely on \emph{autoregressive} (AR) pretrained language models (PLMs) such as GPT~\cite{10.1007/s11023-020-09548-1}, T5~\cite{JMLR:v21:20-074}, and BART~\cite{lewis-etal-2020-bart}, which generate tokens sequentially by conditioning each next token on the prefix~\cite{doi:10.1021/acs.jcim.1c00600,D1RA03086H,irwin2022chemformer}. In this work, while we retain SMILES as the standard input/output interface, we explicitly project these molecules into a Knowledge Graph (KG) view represented as a serialized sequence of atom-bond-atom triplets to achieve better generalization and reduce dependency on specific SMILES syntax rules (e.g., ring number). 

\noindent While effective, autoregressive models are not always well suited to SMILES-centric molecule--language modeling. 
First, decoding is inherently left-to-right, so early mistakes can propagate and are difficult to correct~\cite{nie2025largelanguagediffusionmodels,arriola2025blockdiffusioninterpolatingautoregressive}. 
Second, SMILES encodes long-range syntactic and chemical dependencies (e.g., ring closures, branching, and valid substructures) that require coordinated decisions across distant positions; locally plausible token choices can therefore cascade into globally invalid molecules. 
Third, many objectives are inherently structure-level such as scaffold preservation that are hard to enforce under strict prefix conditioning. 
These limitations motivate non-autoregressive generation mechanisms that can revise all positions jointly, raising a broader question: \emph{can we build models that generate and explain molecules while explicitly preserving global validity and structural controllability?}

\noindent Diffusion language models offer a promising non-autoregressive alternative by iteratively denoising a corrupted representation, enabling holistic updates over all positions at each step and better coordination of long-range dependencies~\cite{li2022diffusionlmimprovescontrollabletext,nie2025largelanguagediffusionmodels}. However, standard diffusion applies noise uniformly across tokens, which can corrupt chemically critical tokens (e.g., ring indices and branching markers) too early in the process and hinder recoverability.

\noindent To address this, we propose \texttt{BiMol-Diff}, a unified diffusion framework that addresses the paired tasks of \emph{text}$\rightarrow$\emph{SMILES} generation and \emph{SMILES}$\rightarrow$\emph{text} captioning under a common denoising formulation. Unlike standard diffusion, \texttt{BiMol-Diff} employs a token-aware noising strategy that assigns token-dependent noise levels using the per-token training loss as a proxy for recovery difficulty. Intuitively, tokens that are consistently harder to denoise receive a more conservative corruption schedule, improving recoverability. We apply this mechanism consistently across both directions: for \emph{SMILES}$\rightarrow$\emph{text}, it helps preserve semantic tokens during denoising, and for \emph{text}$\rightarrow$\emph{SMILES}, it helps preserve chemically salient SMILES tokens during generation.

\noindent \textbf{Contributions.} We make three primary contributions: (1) a token-aware noising strategy which preserves critical structural semantics during the diffusion process; (2) a unified framework that solves the bidirectional tasks of molecule generation and captioning, utilizing the same token-aware mechanism for both inverse problems; and (3) strong empirical performance across a wide range of metrics on both molecule generation and molecule captioning benchmarks.